\documentclass[10pt,a4paper]{amsart}
\usepackage[margin=19mm]{geometry}
\usepackage{amsmath,amssymb,mathtools}
\usepackage[scr=rsfs,cal=euler]{mathalfa}
\usepackage{xfrac}
\usepackage[unicode,hidelinks,bookmarks=false]{hyperref}
\newcommand{\ie}{i.e.\ }
\newcommand{\cf}{cf.\ }
\newcommand{\resp}{respectively\ }
\newcommand{\Subsection}{\subsection}
\providecommand{\nobreakdash}{\nobreak\hbox{-}}
\setlength{\emergencystretch}{2em}
\multlinegap0pt
\usepackage{bm}
\usepackage{bbm}
\usepackage{dsfont}
\usepackage{xspace}
\usepackage{cancel}
\usepackage{mathtools}
\usepackage{amsthm}
\makeatletter
\@ifundefined{thm}{\newtheorem{thm}{Thm}}{}
\@ifundefined{rem}{\newtheorem{rem}[thm]{Remark}}{}
\makeatother
\newcommand{\DD }{\mathop{}\!\mathrm{D}}
\newcommand{\E}{\mathbb E}
\newcommand{\cL}{\mathcal L}
\newcommand{\dd }{\mathop{}\!\mathrm{d}}
\newcommand{\ddt }{\dd t}
\newcommand{\tr}{\mathrm{trace}}
\title[An optimal experimental design approach to sensor placement ]{An optimal experimental design approach to sensor placement in continuous stochastic filtering}
\author{Sahani Pathiraja, Claudia Schillings, Philipp Wacker}
\date{Source: arXiv:2508.12288v1 (CC BY 4.0)}
\begin{document}
\maketitle

\noindent\emph{Source and licence.} An optimal experimental design approach to sensor placement in continuous stochastic filtering. Version arXiv:2508.12288v1 (CC BY 4.0), distributed under the Creative Commons Attribution 4.0 International licence, as recorded in the arXiv licence metadata for this exact version and shown on the abstract page https://arxiv.org/abs/2508.12288v1. Full rights notice: licenses/arxiv-2508.12288-CC-BY-4.0.txt.\par\smallskip

\noindent\textit{Editorial scope.} Counted unit: the Rem whose source label is \texttt{None} in the article (source0.tex lines 624--630, source label \texttt{None}), delivered together with the article's own complete original proof of it (source0.tex lines 631--695, one \texttt{proof} environment of 65 lines). No mathematics is written, repaired or re-derived: statement and proof are the article's own text line for line. Required context retained uncounted: eq:bwd\_KS (lines 615--618). Every definition, display and label of those lines is retained unchanged. Omitted from this package and not redistributed: 1--614, 619--623, 696--880. Nothing inside a retained excerpt is omitted or altered: every retained body is delivered exactly as the article's lines are written, so no omission receipt of any kind accompanies this package. Citations: the retained excerpts carry no citation command at all, so no bibliography entry of the article is transcribed and no reference is redistributed. Reference adaptation: none was needed and none was made; every reference inside the delivered unit resolves inside it, so no unresolved reference or citation is produced. Printed slips kept literal: no printed word, symbol, hypothesis or number of the counted statement or of its proof was altered. Layout and class: the article's own class is \texttt{article} with packages including bm, enumerate, geometry, xcolor, mathrsfs, amsfonts, multirow, amssymb, dsfont, caption, comment, subcaption, blkarray, amsmath; the wrapper is the lane's pinned \texttt{amsart} editorial wrapper at 10pt on A4, which restates the theorem-like environments the retained text needs and adds the article's own macro definitions that the retained text needs (\texttt{\textbackslash DD}, \texttt{\textbackslash E}, \texttt{\textbackslash cL}, \texttt{\textbackslash dd}, \texttt{\textbackslash ddt}, \texttt{\textbackslash tr}), with extra wrapper packages (bm, bbm, dsfont, xspace, cancel, mathtools). The delivered numbering is the unit's own. Assets: no retained span contains a figure environment, a graphics-inclusion command, a PDF-page inclusion, a table or a TikZ picture; no image file is copied into this package, the package declares no asset dependency and no image file is redistributed with it. Licence: version arXiv:2508.12288v1 of this article is distributed under the Creative Commons Attribution 4.0 International licence, as recorded in the arXiv licence metadata for this exact version and shown on the abstract page https://arxiv.org/abs/2508.12288v1. A full rights notice is delivered at licenses/arxiv-2508.12288-CC-BY-4.0.txt. Characters: every retained excerpt is pure ASCII, so the delivered engine, XeLaTeX with its default Latin Modern text font, reports no missing character.
\begin{align}\label{eq:bwd_KS}
    -\frac{\dd\lambda_t(x)}{\dd t} &= -p_t(x) D_{p_t}u_\mathrm{int}(x,p_t(x)) -\nabla \cdot \left(\lambda_t(x) [f(x) -\Sigma \nabla \log p_t(x)] + \frac{1}{2} \Sigma \nabla \lambda_t(x) \right)\\
    \lambda_T(x) &= p_T(x) D_{p_T}u_\mathrm{final}(x,p_T(x))
\end{align}

\begin{rem}
    Using the log-Zakai (instead of the more common Zakai equation, or even the Kushner-Stratonovich equation) has two main benefits: First, we found the log-Zakai equation to be more suitable in our context since we don't need to numerically enforce non-negativity of $p_t$ by working with the log-density directly. Second, if we use the (non-log) Zakai equation, then the (backwards) adjoint equation becomes a stochastically driven semilinear parabolic PDE (like the one for $p_t$ itself) of form
    \begin{equation*}
        \text{adjoint from non-log Zakai:} \quad-\dd \lambda_t(x) = \cL \lambda_t(x)\dd t + p_t(x)\lambda_t(x) g(x)^\top \Gamma(t)^{-1}\dd Z(t),
    \end{equation*}
    while \eqref{eq:bwd_KS} leads to a smooth solution since there is no direct influence of the stochastic forcing $\dd Z(t)$, only via its influence on the filtering distribution $p_t$.
\end{rem}

\begin{proof}
The Lagrangian\footnote{We drop Lagrangian multipliers for nonnegativity and normalisation for simplicity, keeping in mind that the gradient has to be chosen to respect these constraints.} is then (where in the below, $p$ means the entire trajectory $\{p_t\}_{t \in [0,T]}$),
\begin{align*}
    \mathbf L_T(p) = \E U(p) - \E \int_0^T \int \lambda_t(x) \left(\dd  \log p_t(x) - \frac{\cL^\ast p_t(x)}{p_t(x)} \dd t -  g(x)^\top \Gamma^{-1} \dd Z(t) +\frac{1}{2}\xi(t) \| g\|^2_{\Gamma^{-1}}\dd t \right) \dd x
\end{align*}
This means that
\begin{align*}
    D_\xi \mathbf L(p)[h] &=  \E D_p U(p) [D_\xi p [h]] - \E \int_0^T \int \lambda_t(x) \left( D_\xi (\dd \log p_t(x))[h]  - \frac{\cL^\ast D_\xi p_t(x)[h]}{p_t(x)} + \frac{\cL^\ast p_t(x)}{p_t(x)^{2}} D_\xi p_t(x) \dd t \right)\dd x \\
    & - \E \int_0^T\int \lambda_t(x)  \left( - g(x)^\top \Gamma^{-1} \DD_\xi\dd Z(t)  +\frac{1}{2} h(t)\| g(x)\|^2_{\Gamma^{-1}}dt \right) \dd x \\
     \intertext{We perform stochastic integration by parts with respect to time on the integral $\int_0^T \int \lambda_t(x) D_\xi (\dd\log p_t(x))[h]$ (yielding an additional It\=o correction term containing the quadratic covariation between $\lambda_t$ and $D_\xi\log p_t$), and compute the Fr\'echet derivative of the observation path with respect to $\xi$, leading to}
     &= \E D_p U(p) [D_\xi p [h]] - \E \int_0^T \int \lambda_t(x) \left(   - \frac{\cL^\ast D_\xi p_t(x)[h]}{p_t(x)} + \frac{\cL^\ast p_t(x)}{p_t(x)^{2}} D_\xi p_t(x) \dd t \right)\dd x \\
    & - \E \int_0^T\int\lambda_t(x)  \left( - g(x)^\top \Gamma^{-1} \left(g(X(t))h(t)\dd t + \frac{1}{2}h(t) \xi(t)^{-1/2} \Gamma^{1/2} \dd W_t \right)  +\frac{1}{2} h(t)\| g(x)\|^2_{\Gamma^{-1}}dt \right) \dd x \\
    & -\E \int \left( \lambda_T D_\xi  \log p_T[h] - \int_0^T \dd\lambda_t D_\xi \log p_t(x)[h] dt -\frac{1}{2} d\lambda_t d(D_\xi \log p_t[h])   \right) \dd x\\
    \intertext{We now make the ansatz $d\lambda_t = \phi(x, \xi(t)) dW_t + \psi(x, \xi(t))dt $ with $\phi$ and $\psi$ yet to be determined. At the same time, we know that (from the Zakai equation) $d D_\xi \log p_t [h]= \text{[drift terms]}+ g(x)^\top \Gamma^{-1} \frac{h(t)}{2\xi^{1/2}(t)} \Gamma^{1/2} \dd W_t$, so that we get the following general form of the quadratic covariation term: $d\lambda_t d(D_\xi \log p_t[h]) = \frac{g(x)^\top \Gamma^{-1}(t)\phi(x,\xi(t))h(t) }{2\xi^{1/2}(t)}\dd t$. In addition, we recall that $U(p) = U_\mathrm{final}(p_T) + \int_0^T U_\mathrm{int}(p_s)\dd s = \int u_\mathrm{final}(x,p_T(x)))\dd x +  \int_0^T \int u_\mathrm{int}(x,(p_s(x)))\dd x$, and use the adjoint operator to simplify the expression $-\int \lambda_t(x) \frac{\cL^\star D_\xi p_t(x)[h]}{p_t(x)}\dd t = -\int \cL\left(\frac{\lambda_t}{p_t} \right)(x)D_\xi p_t(x)[h]\dd t$, so}
    &= \E \int D_{p_T}u_\mathrm{final}(x,p_T(x))  D_\xi p_T [h](x)\dd x + \E \int_0^T\int D_{p_t}u_\mathrm{int}(x,p_t(x))[D_\xi p_t[h](x)]\dd x\dd t\\
    &- \E \int_0^T \int \left[ -\mathcal{L} \left( \frac{\lambda_t}{p_t} \right)(x)   
+\frac{\lambda_t(x)}{p_t(x)} \left(     \frac{\cL^\ast p_t(x)}{p_t(x)}  \right)\right]D_\xi p_t(x)  \dd t \\
    & - \E \int_0^T\int\lambda_t(x)  \left( - g(x)^\top \Gamma^{-1} \left(g(X(t))h(t)\dd t + \frac{1}{2}h(t) \xi(t)^{-1/2} \Gamma^{1/2} \dd W_t \right)  +\frac{1}{2} h(t)\| g(x)\|^2_{\Gamma^{-1}}dt \right) \dd x \\
    & - \E \int\left( \lambda_T \frac{D_\xi   p_T[h]}{p_T(x)} - \int_0^T \dd\lambda_t(x) \frac{D_\xi p_t(x)[h]}{p_t(x)}  -\frac{1}{4} \int_0^T \frac{h(t)}{\xi^{1/2} }\Gamma^{1/2} \phi dt   \right) \dd x\\
    &= \E \int \left(D_{p_T}u_\mathrm{final}(x,p_T(x)) - \frac{\lambda_T(x)}{p_T(x)}\right)  D_\xi p_T [h](x)\dd x\\
    &+\E \int_0^T \int D_\xi p_t[h](x) \left(D_{p_t}u_\mathrm{int}(x,p_t(x))\ddt + \cL \left(\frac{\lambda_t}{p_t} \right)(x)\ddt - \frac{\lambda_t(x)}{p_t(x)}\frac{\cL^\star p_t(x)}{p_t(x)}\ddt+ \frac{\dd \lambda_t(x)}{p_t(x)}\right)\dd x\\
    &+ \E \int_0^Th(t)\int \lambda_t(x)  \left( -\langle g(x), g(X(t))\rangle_{\Gamma(t)} + \frac{1}{2}\|g(x)\|_\Gamma(t)^2 \right)   + \frac{1}{4\xi^{1/2}(t)}\Gamma^{1/2}(t)\phi(x,\xi)\dd t \dd x\\
    &+\E \int_0^T\int \lambda_t(x)h(t)\left( - \frac{1}{2}g(x)^\top \Gamma^{-1/2} \xi(t)^{-1/2}  \dd W_t + \frac{1}{4\xi^{1/2}(t)}\Gamma^{1/2}(t)\phi(x,\xi(t))\right)
\end{align*}
The four lines in the last expression determine the following components of interest:

By setting $\lambda_T(x) = p_T(x) D_{p_T}u_\mathrm{final}(x,p_T(x))$, the first line vanishes. The second line can be dropped by posing the (backwards) adjoint differential equation
\begin{align*}
    -\partial_t\lambda_t(x) = -\lambda_t(x) \frac{(\cL^\star p_t)(x)}{p_t(x)} + \cL\left( \frac{\lambda_t(x)}{p_t(x)}\right)(x) + D_{p_t}u_\mathrm{int}(x,p_t(x))
\end{align*}
With this choice, we see that the diffusion coefficient of the dynamics of $\lambda$ is $0$, so $\phi = 0$, and the term stemming from the quadratic covariation vanishes. Since we take the expectation with respect to $W_t$,  the stochastic integral is dropped, as well (under mild integrability assumptions on the integrand, making sure the local martingale given by the It\=o integral is actually a martingale). This means that the fourth line vanishes. The third line, finally, evaluates to an expression yielding the Fr\'echet derivative of the utility functional with respect to $\xi$:


\[\mathbb{E}D_\xi \mathbf L(p) = \E \int   \lambda_t(x)  \left( - g(x)^\top \Gamma^{-1} g(X(t))    +\frac{1}{2} \| g(x)\|^2_{\Gamma^{-1}} \right) \dd x\]


 We now put in a bit more effort to simplify the first two terms in the adjoint equation, using the special structure of $\cL$ in our case where 
\begin{align*}
    \cL^\star p_t(x) &= -\nabla \cdot (p_t(x) f(x)) + \frac{1}{2} \tr[\Sigma \nabla^2 p_t(x)] \\
    \cL p_t(x)  & = f(x) \cdot \nabla p_t(x) + \frac{1}{2}\tr[\Sigma \nabla^2 p_t(x)]
\end{align*}
We can then define $\mu_t = \frac{\lambda_t}{p_t}$ and evaluate the term in the right hand side of the adjoint equation as
\begin{align*}
    - p_t(x)\cL \left(\frac{\lambda_t}{p_t} \right)(x) + \lambda_t(x)\frac{\cL^\star p_t(x)}{p_t(x)} &= -p_t \left(  f \cdot \nabla \left(\frac{\lambda_t}{p_t} \right) + \frac{1}{2}\tr\left[\Sigma \nabla^2 \left( \frac{\lambda_t}{p_t} \right) \right] \right) \\
    &+  \frac{\lambda_t}{p_t} \left( -\nabla \cdot (p_t(x) f(x)) + \frac{1}{2} \tr[\Sigma \nabla^2 p_t(x)] \right)  \\
    & = -\nabla \cdot \left( \lambda_t f \right) + \frac{1}{2} \tr \left[ \frac{\lambda_t}{p_t} \nabla^2 p_t \Sigma - p_t \nabla^2 \left( \frac{\lambda_t}{p_t} \right) \Sigma  \right] \\
    &= -\nabla \cdot \left( \lambda_t f \right)  + \frac{1}{2} \tr\left[ \Sigma \left( p \nabla^2 \mu - \mu \nabla^2 p \right) \right].
\end{align*}


We now simplify the second-order term using the identity:
\[
p \nabla^2 \mu - \mu \nabla^2 p = \nabla \cdot (p \nabla \mu - \mu \nabla p)
\]

Hence $\tr\left[ \Sigma (p \nabla^2 \mu - \mu \nabla^2 p) \right]
= \nabla \cdot \left( \Sigma (p \nabla \mu - \mu \nabla p) \right) $. 

So the adjoint equation becomes ($\mu = \lambda /p$) is 
\begin{align*}
     \partial_t \lambda_t(x) &= -p_t(x) D_{p_t}u_\mathrm{int}(x,p_t(x)) -\nabla \cdot \left(\lambda_t(x) f(x) + \frac{1}{2} \Sigma \left(p_t(x) \nabla \frac{\lambda_t(x)}{p_t(x)} - \frac{\lambda_t(x)}{p_t(x)} \nabla p_t(x)\right) \right)\\
     \intertext{then we use $p_t(x) \nabla \frac{\lambda_t(x)}{p_t(x)} = \nabla \lambda_t(x) - \lambda_t(x) \nabla \log p_t(x)$, so we get}
     &= -p_t(x) D_{p_t}u_\mathrm{int}(x,p_t(x)) -\nabla \cdot \left(\lambda_t(x) [f(x) -\Sigma \nabla \log p_t(x)] + \frac{1}{2} \Sigma \nabla \lambda_t(x) \right)
\end{align*}
\end{proof}

\end{document}
